\begin{table}[!ht]
\centering
\caption{H2 --- Wage--price connectedness is stronger when inflation is high}
\label{tab:h2}
\setlength{\tabcolsep}{6pt}\footnotesize
\begin{threeparttable}
\begin{tabular}{l c c c}
\multicolumn{4}{l}{\textit{Panel A. Two-variable wage--CPI system (monthly, full sample), by quantile}}\\
\toprule
 & $\tau=0.5$ & $\tau=0.7$ & $\tau=0.9$ \\
\midrule
Wage--CPI total connectedness (TCI) & 10.8 & 16.6 & 21.6 \\
\quad lagged Wages$\to$CPI & 5.0 & 7.7 & 11.5 \\
\quad lagged CPI$\to$Wages & 6.2 & 9.2 & 16.1 \\
\bottomrule
\end{tabular}

\vspace{4pt}
\begin{tabular}{l c c c}
\multicolumn{4}{l}{\textit{Panel B. By subperiod (monthly)}}\\
\toprule
 & $\tau=0.5$ & $\tau=0.7$ & $\tau=0.9$ \\
\midrule
\multicolumn{4}{l}{\textit{\;1964--2019 (with Great Inflation)}}\\
Wage--CPI total connectedness (TCI) & 13.3 & 18.4 & 25.2 \\
\quad lagged Wages$\to$CPI & 6.0 & 9.1 & 15.5 \\
\quad lagged CPI$\to$Wages & 5.9 & 9.0 & 14.9 \\
\addlinespace
\multicolumn{4}{l}{\textit{\;1983--2026 (post-Great-Inflation)}}\\
Wage--CPI total connectedness (TCI) & 1.2 & 1.8 & 3.6 \\
\quad lagged Wages$\to$CPI & 0.2 & 0.2 & 1.3 \\
\quad lagged CPI$\to$Wages & 0.9 & 1.7 & 3.0 \\
\bottomrule
\end{tabular}

\vspace{4pt}
\begin{tabular}{l c c c}
\multicolumn{4}{l}{\textit{Panel C. Quarterly (full sample), by quantile}}\\
\toprule
 & $\tau=0.5$ & $\tau=0.7$ & $\tau=0.9$ \\
\midrule
Wage--CPI total connectedness (TCI) & 25.7 & 32.6 & 53.4 \\
\quad lagged Wages$\to$CPI & 14.6 & 19.3 & 35.1 \\
\quad lagged CPI$\to$Wages & 13.0 & 18.5 & 23.8 \\
\bottomrule
\end{tabular}
\begin{tablenotes}[flushleft]\footnotesize
\item \textit{Notes:} Two-variable system of wage growth (AHETPI) and overall CPI (CPIAUCSL); pseudo-quantile $R^2$ connectedness (\%), Genizi estimator, $n_{\text{lag}}=2$. Wage--price connectedness rises sharply into the high-inflation tail at all frequencies, and is far stronger in the Great-Inflation-inclusive subperiod (1964--2019) than post-1983. In the high tail the lagged components dominate. Quarterly ($n_{\text{lag}}=2$) spans a six-month horizon.
\end{tablenotes}
\end{threeparttable}
\end{table}

